\documentclass[11pt,a4paper]{article}

% Packages
\usepackage[utf8]{inputenc}
\usepackage[T1]{fontenc}
\usepackage{amsmath}
\usepackage{amsfonts}
\usepackage{amssymb}
\usepackage{graphicx}
\usepackage{booktabs}
\usepackage{longtable}
\usepackage{array}
\usepackage{multirow}
\usepackage{wrapfig}
\usepackage{float}
\usepackage{xcolor}
\usepackage{geometry}
\usepackage{enumitem}
\usepackage{hyperref}
\usepackage{fancyhdr}
\usepackage{titlesec}
\usepackage{pgfplots}
\usepackage{tikz}
\usepackage{url}
\usepackage{cite}
\usepackage{listings}
\usepackage{caption}
\usepackage{subcaption}

% Page setup
\geometry{margin=1in}
\pagestyle{fancy}
\fancyhf{}
\fancyhead[LE,RO]{\thepage}
\fancyhead[LO,RE]{\leftmark}

% Custom colors
\definecolor{mlcolor}{RGB}{52, 152, 219}
\definecolor{statcolor}{RGB}{231, 76, 60}
\definecolor{accentcolor}{RGB}{46, 204, 113}
\definecolor{dlcolor}{RGB}{255, 165, 0}

% Hyperref setup
\hypersetup{
    colorlinks=true,
    linkcolor=blue,
    filecolor=magenta,      
    urlcolor=cyan,
    citecolor=green
}

% Title information
\title{Bridging the Gap: Statistical vs Machine Learning Methods\\
\large Performance, Cost, and Interpretability Trade-offs\\
\large in the Age of Transparent AI}
\author{Abderrahmane Moujar \and Oudoum Ali Houmed \and Said Abolhassan Razavi\\
\small Master Paris-Saclay, Applied Statistics}
\date{\today}

\begin{document}

\maketitle

\begin{abstract}
This comprehensive study presents a systematic comparison between statistical and machine learning (ML) approaches across multiple problem domains. We evaluate 13 model-dataset combinations spanning four distinct problem types: regression, classification, time series analysis, and factorial analysis. Our analysis quantifies the fundamental trade-offs between three critical dimensions: \textcolor{mlcolor}{performance} (accuracy, R², RMSE), \textcolor{statcolor}{cost} (training time, computational resources), and \textcolor{accentcolor}{interpretability} (explainability, statistical inference). 

Our findings reveal that statistical models excel in interpretability (88.3/100 vs 37.4/100) and efficiency (20-100x faster training), while ML/DL models demonstrate superior performance on complex datasets (0.818 vs 0.547 average). The study connects these findings to the emerging field of Interpretable ML and Transparent AI, providing evidence-based recommendations for model selection. We demonstrate that no single approach dominates all categories, and the optimal choice depends on specific requirements including regulatory constraints, dataset characteristics, and computational resources.

\textbf{Keywords:} Machine Learning, Statistical Models, Interpretability, Transparent AI, Model Comparison, Performance Trade-offs
\end{abstract}

\tableofcontents
\newpage

\section{Introduction}

\subsection{Background and Motivation}

The landscape of predictive modeling has been fundamentally shaped by two distinct yet interconnected paradigms: traditional statistical methods and modern machine learning approaches. While statistical models have provided the foundation for scientific inference and hypothesis testing for over a century, machine learning methods have revolutionized predictive accuracy in complex, high-dimensional problems.

However, as artificial intelligence systems increasingly inform critical decisions in healthcare, finance, and public policy, the demand for \textbf{transparent and interpretable} models has reached unprecedented levels. This convergence of performance requirements and interpretability needs has created a fundamental tension: how can practitioners balance the superior predictive power of complex ML models with the transparency and trust required for high-stakes applications?

The emergence of \textbf{Interpretable ML} and \textbf{Transparent AI} research represents a concerted effort to bridge this gap. As noted in recent work on transparent AI \cite{transparent_ai_2025}, the challenge is not merely technical but encompasses organizational, regulatory, and human factors. This study contributes to this ongoing dialogue by providing empirical evidence on the trade-offs between statistical and ML approaches.

\subsection{Problem Statement}

The core research question driving this investigation is:

\textit{How do Machine Learning and Statistical models compare across different problem domains in terms of:}
\begin{itemize}
    \item \textcolor{mlcolor}{\textbf{Performance}} (accuracy, R², RMSE)
    \item \textcolor{statcolor}{\textbf{Cost}} (training time, computational resources)
    \item \textcolor{accentcolor}{\textbf{Interpretability}} (explainability, statistical inference)
\end{itemize}

This question is particularly relevant in the context of regulatory requirements (EU AI Act, Canada AIDA, US AI Bill of Rights) that mandate transparency for high-risk AI systems. Understanding these trade-offs enables evidence-based model selection that aligns with both performance objectives and compliance requirements.

\subsection{Research Objectives}

This study aims to:

\begin{enumerate}
    \item \textbf{Systematic Comparison:} Compare ML and Statistical approaches across multiple problem domains using standardized evaluation protocols
    \item \textbf{Quantify Trade-offs:} Provide quantitative measures of the performance-cost-interpretability trade-offs
    \item \textbf{Evidence-based Recommendations:} Develop practical guidelines for model selection based on empirical findings
    \item \textbf{Connect to Transparent AI:} Link findings to the broader research agenda on interpretable ML and transparent AI systems
\end{enumerate}

\subsection{Scope of the Study}

Our comprehensive analysis encompasses:

\begin{table}[h]
\centering
\begin{tabular}{lrr}
\toprule
\textbf{Category} & \textbf{Count} & \textbf{Details} \\
\midrule
Problem Domains & 4 & Regression, Classification, Time Series, Factorial Analysis \\
Datasets & 13 & 150 to 35,000 samples \\
ML Models & 5 & Linear Regression, Random Forest, SVM, RNN/LSTM, Gradient Boosting \\
Statistical Models & 5 & OLS Regression, Logistic Regression, ANOVA, ARIMA, Cox/Poisson \\
Total Comparisons & 13 & Model-Dataset pairs \\
\bottomrule
\end{tabular}
\caption{Study Scope Overview}
\label{tab:scope}
\end{table}

\subsection{Problem Domains and Datasets}

\subsubsection{Regression Analysis}
\begin{itemize}
    \item \textbf{California Housing:} 20,640 samples, 8 features - Predicting median house values
    \item \textbf{Diabetes:} 442 samples, 10 features - Disease progression prediction
    \item \textbf{Wine Quality:} 1,599 samples, 11 features - Quality score prediction
    \item \textbf{Energy Consumption:} 1,000 samples (synthetic), 6 features - Energy usage forecasting
\end{itemize}

\subsubsection{Classification}
\begin{itemize}
    \item \textbf{Iris:} 150 samples, 3 classes, 4 features - Classic classification benchmark
    \item \textbf{Wine Classification:} 178 samples, 3 classes, 13 features - Wine type classification
    \item \textbf{Breast Cancer:} 569 samples, 2 classes, 30 features - Medical diagnosis
    \item \textbf{Credit Default:} 1,000 samples (synthetic), 2 classes, 10 features - Financial risk assessment
\end{itemize}

\subsubsection{Time Series Analysis}
\begin{itemize}
    \item \textbf{Energy Consumption:} 1,000 time steps (synthetic) - Temporal energy patterns
    \item \textbf{Temperature:} 1,000 time steps (synthetic) - Climate data forecasting
    \item \textbf{Sales Forecasting:} 1,000 time steps (synthetic) - Business metrics prediction
\end{itemize}

\subsubsection{Factorial Analysis}
\begin{itemize}
    \item \textbf{Treatment Effects:} Clinical trial data - Medical intervention analysis
    \item \textbf{Educational Outcomes:} Teaching method comparison - Educational research
    \item \textbf{Agricultural Experiment:} Fertilizer × Water level (3×3 factorial) - Agricultural science
    \item \textbf{Drug Efficacy Study:} Drug treatment comparison - Pharmaceutical research
\end{itemize}

\section{Literature Review and Theoretical Framework}

\subsection{The Interpretability-Performance Trade-off}

The fundamental tension between model performance and interpretability has been extensively documented in the machine learning literature. This trade-off can be mathematically expressed as:

\begin{equation}
\max_\theta \quad P(\theta) - \alpha C(\theta) + \beta I(\theta)
\label{eq:optimization}
\end{equation}

where:
\begin{itemize}
    \item $P(\theta)$ represents \textbf{Performance} (e.g., R², accuracy)
    \item $C(\theta)$ represents \textbf{Cost} (training time, computational resources)
    \item $I(\theta)$ represents \textbf{Interpretability} (explainability score $\in [0, 100]$)
    \item $\alpha$ and $\beta$ are weighting factors reflecting domain-specific priorities
\end{itemize}

Figure~\ref{fig:interpretability_tradeoff} illustrates this relationship, showing how different model types fall along the performance-interpretability spectrum.

\begin{figure}[h]
\centering
\includegraphics[width=0.8\textwidth]{bridging_gap.png}
\caption{The Performance-Interpretability Trade-off. The purple curve represents the traditional inverse relationship, while the red arrow indicates the direction of Interpretable ML research toward the ideal solution (high performance, high interpretability).}
\label{fig:interpretability_tradeoff}
\end{figure}

\subsection{Statistical Models: Foundations and Characteristics}

Statistical models are grounded in mathematical theory and provide:

\begin{itemize}
    \item \textbf{Statistical Inference:} P-values, confidence intervals, hypothesis testing
    \item \textbf{Model Diagnostics:} F-tests, AIC, BIC, residual analysis
    \item \textbf{Theoretical Guarantees:} Asymptotic properties, convergence results
    \item \textbf{Interpretability:} Direct coefficient interpretation, causal inference capabilities
\end{itemize}

\subsubsection{Ordinary Least Squares (OLS) Regression}

The OLS regression model is defined as:

\begin{equation}
Y = X\beta + \varepsilon
\end{equation}

with the least squares estimator:

\begin{equation}
\hat{\beta} = (X^T X)^{-1} X^T Y
\end{equation}

This approach provides:
\begin{itemize}
    \item Unbiased estimates under Gauss-Markov assumptions
    \item Statistical significance testing via t-tests and F-tests
    \item Confidence intervals for coefficients
    \item Model diagnostics (R², adjusted R², residual analysis)
\end{itemize}

\subsubsection{Logistic Regression}

For binary classification, logistic regression models:

\begin{equation}
P(Y=1|X) = \frac{1}{1 + e^{-X\beta}}
\end{equation}

with the log-odds transformation:

\begin{equation}
\log\left(\frac{p}{1-p}\right) = X\beta
\end{equation}

This provides odds ratios, p-values, and confidence intervals for each coefficient.

\subsubsection{ANOVA (Analysis of Variance)}

ANOVA tests for differences between group means using the model:

\begin{equation}
Y_{ij} = \mu + \alpha_i + \beta_j + (\alpha\beta)_{ij} + \varepsilon_{ij}
\end{equation}

with the F-statistic:

\begin{equation}
F = \frac{\text{MS}_{\text{between}}}{\text{MS}_{\text{within}}}
\end{equation}

\subsubsection{ARIMA (AutoRegressive Integrated Moving Average)}

For time series, ARIMA models capture temporal dependencies:

\begin{equation}
\phi(B)(1-B)^d X_t = \theta(B)\varepsilon_t
\end{equation}

where $(p,d,q)$ parameters are selected via information criteria.

\subsection{Machine Learning Models: Capabilities and Limitations}

Modern ML models excel at:

\begin{itemize}
    \item \textbf{Complex Pattern Recognition:} Non-linear relationships, interactions
    \item \textbf{High-Dimensional Data:} Feature selection, dimensionality reduction
    \item \textbf{Scalability:} Large datasets, distributed computing
    \item \textbf{Adaptive Learning:} Online learning, transfer learning
\end{itemize}

However, they typically lack:

\begin{itemize}
    \item Direct statistical inference capabilities
    \item Theoretical guarantees (except in specific cases)
    \item Built-in interpretability mechanisms
    \item Hypothesis testing frameworks
\end{itemize}

\subsubsection{Random Forest}

Random Forest combines multiple decision trees:

\begin{equation}
\hat{y} = \frac{1}{B}\sum_{b=1}^B T_b(X; \theta_b)
\end{equation}

where $T_b \sim \text{Decision Tree}$ and $B$ is the number of trees.

\subsubsection{Support Vector Machines (SVM)}

SVM solves the optimization problem:

\begin{equation}
\min_{\mathbf{w},b} \quad \frac{1}{2}\|\mathbf{w}\|^2 + C\sum_{i=1}^n \xi_i
\end{equation}

with kernel function:

\begin{equation}
K(\mathbf{x}_i, \mathbf{x}_j) = \phi(\mathbf{x}_i)^T\phi(\mathbf{x}_j)
\end{equation}

\subsubsection{Recurrent Neural Networks (LSTM)}

LSTM networks model sequential dependencies:

\begin{equation}
h_t = \tanh(W_h h_{t-1} + W_x x_t + b)
\end{equation}

\begin{equation}
\hat{y}_t = f(h_t)
\end{equation}

\subsection{Transparent AI and Interpretable ML}

Recent research, as exemplified by the work on Transparent AI \cite{transparent_ai_2025}, emphasizes that interpretability is not optional for high-stakes AI deployment. Key principles include:

\begin{enumerate}
    \item \textbf{Early Integration:} Interpretability must be embedded from initial design, not retrofitted
    \item \textbf{Stakeholder-Specific Needs:} Different audiences require different explanation types
    \item \textbf{Comprehensive Evaluation:} Beyond accuracy metrics, evaluate explanation quality
    \item \textbf{Regulatory Alignment:} Ensure compliance with emerging AI governance frameworks
\end{enumerate}

\section{Methodology}

\subsection{Experimental Design}

Our experimental design follows a systematic comparison framework ensuring fair evaluation across all model-dataset combinations.

\subsection{Data Preprocessing}

\subsubsection{Exploratory Data Analysis (EDA)}

For each dataset, we performed comprehensive EDA:

\begin{itemize}
    \item \textbf{Distribution Analysis:} Histograms, box plots, Q-Q plots
    \item \textbf{Correlation Analysis:} Correlation matrices, multicollinearity detection
    \item \textbf{Missing Value Detection:} Identification and treatment strategies
    \item \textbf{Outlier Identification:} Statistical tests (IQR, Z-score) and visualization
    \item \textbf{Feature-Target Relationships:} Scatter plots, partial dependence plots
\end{itemize}

\subsubsection{Preprocessing Pipeline}

Standardized preprocessing steps applied consistently:

\begin{enumerate}
    \item \textbf{Train-Test Split:} 80-20 split, maintaining temporal order for time series
    \item \textbf{Standardization:} $z = \frac{x - \mu}{\sigma}$ for numerical features
    \item \textbf{Feature Scaling:} Min-max scaling for numerical stability
    \item \textbf{Handling Skewed Distributions:} Log transformation, Box-Cox transformation
    \item \textbf{Outlier Treatment:} Winsorization or removal based on domain knowledge
    \item \textbf{Categorical Encoding:} One-hot encoding for nominal, ordinal encoding for ordinal
\end{enumerate}

\textbf{Critical Consistency Requirement:} Both statistical and ML models were trained on identical preprocessed data to ensure fair comparison.

\subsection{Model Implementation}

\subsubsection{Statistical Models}

\begin{itemize}
    \item \textbf{OLS Regression:} Implemented using \texttt{statsmodels} (Python)
    \item \textbf{Logistic Regression:} \texttt{statsmodels} with maximum likelihood estimation
    \item \textbf{ANOVA:} \texttt{scipy.stats} and \texttt{statsmodels}
    \item \textbf{ARIMA:} \texttt{statsmodels.tsa.arima.model} with auto-selection of $(p,d,q)$
    \item \textbf{Cox/Poisson Regression:} \texttt{lifelines} and \texttt{statsmodels}
\end{itemize}

\subsubsection{Machine Learning Models}

\begin{itemize}
    \item \textbf{Linear Regression:} \texttt{scikit-learn} with default hyperparameters
    \item \textbf{Random Forest:} \texttt{scikit-learn} with grid search for hyperparameter tuning
    \item \textbf{SVM:} \texttt{scikit-learn} with RBF kernel and cross-validation
    \item \textbf{RNN/LSTM:} \texttt{PyTorch} implementation with appropriate architecture
    \item \textbf{Gradient Boosting:} \texttt{scikit-learn} GradientBoostingRegressor/Classifier
\end{itemize}

\subsection{Evaluation Metrics}

\subsubsection{Performance Metrics}

\begin{itemize}
    \item \textbf{Regression:} R² (coefficient of determination), RMSE (Root Mean Squared Error), MAE (Mean Absolute Error)
    \item \textbf{Classification:} Accuracy, Precision, Recall, F1-score, AUC-ROC
    \item \textbf{Time Series:} R², RMSE, MAE on test set forecasts
\end{itemize}

\subsubsection{Cost Metrics}

\begin{itemize}
    \item \textbf{Training Time:} Wall-clock time in seconds
    \item \textbf{Prediction Time:} Inference time per sample
    \item \textbf{Memory Usage:} Peak memory consumption during training
    \item \textbf{Efficiency Score:} $\frac{\text{Performance}}{\text{Cost}}$ (normalized)
\end{itemize}

\subsubsection{Interpretability and Transparency Metrics}

We developed a comprehensive composite interpretability and transparency score (0-100) that quantifies how well a model can be understood, explained, and trusted by stakeholders. The score is computed as a weighted sum of four key dimensions:

\begin{equation}
I(\theta) = I_{\text{transparency}}(\theta) + I_{\text{inference}}(\theta) + I_{\text{explanation}}(\theta) + I_{\text{documentation}}(\theta)
\label{eq:interpretability}
\end{equation}

where $I(\theta) \in [0, 100]$ and each component is scored independently.

\paragraph{Component 1: Model Transparency (0-40 points)}

Model transparency measures the extent to which model parameters and structure can be directly understood without additional explanation methods.

\textbf{Scoring Criteria:}
\begin{itemize}
    \item \textbf{40 points:} Fully transparent models where every parameter has direct, intuitive meaning
    \begin{itemize}
        \item Examples: OLS Regression (coefficients = unit change in outcome per unit change in predictor)
        \item Logistic Regression (coefficients = log-odds ratios)
        \item Decision Trees (if-then rules directly interpretable)
    \end{itemize}
    
    \item \textbf{30 points:} Mostly transparent with some complexity
    \begin{itemize}
        \item Examples: Generalized Additive Models (GAMs) - smooth functions interpretable
        \item ANOVA - group means and effects interpretable
    \end{itemize}
    
    \item \textbf{20 points:} Partially transparent, requires domain knowledge
    \begin{itemize}
        \item Examples: ARIMA models - coefficients have statistical meaning but require time series expertise
        \item Cox Regression - hazard ratios interpretable with survival analysis knowledge
    \end{itemize}
    
    \item \textbf{10 points:} Low transparency, complex interactions
    \begin{itemize}
        \item Examples: Random Forest - individual trees interpretable, ensemble less so
        \item Gradient Boosting - sequential improvements hard to interpret
    \end{itemize}
    
    \item \textbf{0 points:} Black box models
    \begin{itemize}
        \item Examples: Deep Neural Networks - weights have no direct meaning
        \item Complex ensemble methods with non-linear interactions
    \end{itemize}
\end{itemize}

\textbf{Computation:}
\begin{equation}
I_{\text{transparency}}(\theta) = \begin{cases}
40 & \text{if fully transparent (linear models, simple trees)} \\
30 & \text{if mostly transparent (GAMs, ANOVA)} \\
20 & \text{if partially transparent (ARIMA, Cox)} \\
10 & \text{if low transparency (ensembles)} \\
0 & \text{if black box (deep learning)}
\end{cases}
\end{equation}

\paragraph{Component 2: Statistical Inference (0-30 points)}

Statistical inference measures the availability and quality of statistical tools for hypothesis testing, uncertainty quantification, and model validation.

\textbf{Scoring Criteria:}
\begin{itemize}
    \item \textbf{30 points:} Full statistical inference suite
    \begin{itemize}
        \item P-values for all coefficients/parameters
        \item Confidence intervals (95\% CI standard)
        \item Hypothesis tests (t-tests, F-tests, likelihood ratio tests)
        \item Model diagnostics (residual analysis, goodness-of-fit tests)
        \item Examples: OLS Regression, Logistic Regression, ANOVA, ARIMA
    \end{itemize}
    
    \item \textbf{20 points:} Partial statistical inference
    \begin{itemize}
        \item Some parameters have p-values or confidence intervals
        \item Limited hypothesis testing capabilities
        \item Examples: Cox Regression (hazard ratios with CI), Poisson Regression
    \end{itemize}
    
    \item \textbf{10 points:} Minimal statistical inference
    \begin{itemize}
        \item Feature importance scores available
        \item Some uncertainty quantification (e.g., prediction intervals)
        \item Examples: Random Forest (feature importance), Gradient Boosting
    \end{itemize}
    
    \item \textbf{0 points:} No statistical inference
    \begin{itemize}
        \item No p-values, confidence intervals, or hypothesis tests
        \item Examples: Deep Neural Networks, complex ML models without statistical framework
    \end{itemize}
\end{itemize}

\textbf{Computation:}
\begin{equation}
I_{\text{inference}}(\theta) = \begin{cases}
30 & \text{if full inference (p-values, CI, tests, diagnostics)} \\
20 & \text{if partial inference (some CI, limited tests)} \\
10 & \text{if minimal inference (feature importance only)} \\
0 & \text{if no inference}
\end{cases}
\end{equation}

\paragraph{Component 3: Explanation Methods Compatibility (0-20 points)}

This component evaluates how well the model can be explained using post-hoc explanation methods, even if the model itself is not inherently interpretable.

\textbf{Scoring Criteria:}
\begin{itemize}
    \item \textbf{20 points:} Excellent compatibility with multiple explanation methods
    \begin{itemize}
        \item Works well with SHAP (SHapley Additive exPlanations)
        \item Compatible with LIME (Local Interpretable Model-agnostic Explanations)
        \item Supports partial dependence plots
        \item Feature importance methods applicable
        \item Examples: Random Forest, Gradient Boosting, Linear Models
    \end{itemize}
    
    \item \textbf{15 points:} Good compatibility with some methods
    \begin{itemize}
        \item Works with some explanation methods but not all
        \item May require model-specific adaptations
        \item Examples: SVM (works with SHAP, limited LIME compatibility)
    \end{itemize}
    
    \item \textbf{10 points:} Limited compatibility
    \begin{itemize}
        \item Only basic explanation methods work
        \item Explanations may be unstable or incomplete
        \item Examples: Some deep learning models
    \end{itemize}
    
    \item \textbf{5 points:} Very limited compatibility
    \begin{itemize}
        \item Few explanation methods applicable
        \item Explanations require significant technical expertise
        \item Examples: Complex RNN/LSTM architectures
    \end{itemize}
    
    \item \textbf{0 points:} No effective explanation methods
    \begin{itemize}
        \item Standard explanation methods fail or provide unreliable results
        \item Examples: Extremely complex or proprietary models
    \end{itemize}
\end{itemize}

\textbf{Computation:}
\begin{equation}
I_{\text{explanation}}(\theta) = \begin{cases}
20 & \text{if excellent (SHAP, LIME, PDP all work well)} \\
15 & \text{if good (some methods work)} \\
10 & \text{if limited (basic methods only)} \\
5 & \text{if very limited (few methods)} \\
0 & \text{if none (no effective methods)}
\end{cases}
\end{equation}

\paragraph{Component 4: Documentation and Theoretical Foundation (0-10 points)}

This component assesses the quality of model documentation, theoretical understanding, and educational resources available.

\textbf{Scoring Criteria:}
\begin{itemize}
    \item \textbf{10 points:} Excellent documentation and strong theoretical foundation
    \begin{itemize}
        \item Well-established theoretical framework (decades of research)
        \item Comprehensive documentation with examples
        \item Extensive educational resources (textbooks, courses)
        \item Clear mathematical derivations available
        \item Examples: OLS Regression, Logistic Regression, ANOVA
    \end{itemize}
    
    \item \textbf{7 points:} Good documentation and theoretical understanding
    \begin{itemize}
        \item Solid theoretical foundation
        \item Good documentation
        \item Some educational resources available
        \item Examples: ARIMA, Cox Regression, Random Forest
    \end{itemize}
    
    \item \textbf{5 points:} Moderate documentation
    \begin{itemize}
        \item Some theoretical understanding
        \item Basic documentation available
        \item Limited educational resources
        \item Examples: Gradient Boosting, SVM
    \end{itemize}
    
    \item \textbf{3 points:} Limited documentation
    \begin{itemize}
        \item Emerging or complex methods
        \item Documentation may be technical or incomplete
        \item Examples: Some deep learning architectures
    \end{itemize}
    
    \item \textbf{0 points:} Poor or no documentation
    \begin{itemize}
        \item New or proprietary methods
        \item Minimal documentation
        \item Examples: Cutting-edge or experimental models
    \end{itemize}
\end{itemize}

\textbf{Computation:}
\begin{equation}
I_{\text{documentation}}(\theta) = \begin{cases}
10 & \text{if excellent (strong theory, comprehensive docs)} \\
7 & \text{if good (solid theory, good docs)} \\
5 & \text{if moderate (some theory, basic docs)} \\
3 & \text{if limited (emerging, technical docs)} \\
0 & \text{if poor (minimal or no docs)}
\end{cases}
\end{equation}

\paragraph{Total Interpretability Score Calculation}

The final interpretability and transparency score is computed as:

\begin{equation}
I_{\text{total}}(\theta) = I_{\text{transparency}}(\theta) + I_{\text{inference}}(\theta) + I_{\text{explanation}}(\theta) + I_{\text{documentation}}(\theta)
\end{equation}

where $I_{\text{total}}(\theta) \in [0, 100]$.

\paragraph{Example Calculations}

\textbf{Example 1: OLS Regression}
\begin{align*}
I_{\text{transparency}} &= 40 \quad \text{(coefficients directly interpretable)} \\
I_{\text{inference}} &= 30 \quad \text{(full p-values, CI, tests)} \\
I_{\text{explanation}} &= 20 \quad \text{(excellent SHAP, LIME compatibility)} \\
I_{\text{documentation}} &= 10 \quad \text{(excellent theory and documentation)} \\
I_{\text{total}} &= 40 + 30 + 20 + 10 = \mathbf{100/100}
\end{align*}

\textbf{Example 2: Random Forest}
\begin{align*}
I_{\text{transparency}} &= 10 \quad \text{(ensemble complexity)} \\
I_{\text{inference}} &= 10 \quad \text{(feature importance only)} \\
I_{\text{explanation}} &= 20 \quad \text{(excellent SHAP, LIME compatibility)} \\
I_{\text{documentation}} &= 7 \quad \text{(good documentation)} \\
I_{\text{total}} &= 10 + 10 + 20 + 7 = \mathbf{47/100}
\end{align*}

\textbf{Example 3: Deep Neural Network (LSTM)}
\begin{align*}
I_{\text{transparency}} &= 0 \quad \text{(black box)} \\
I_{\text{inference}} &= 0 \quad \text{(no statistical inference)} \\
I_{\text{explanation}} &= 5 \quad \text{(very limited explanation methods)} \\
I_{\text{documentation}} &= 5 \quad \text{(moderate documentation)} \\
I_{\text{total}} &= 0 + 0 + 5 + 5 = \mathbf{10/100}
\end{align*}

\paragraph{Transparency vs Interpretability}

While often used interchangeably, we distinguish:
\begin{itemize}
    \item \textbf{Transparency:} The model's internal structure and parameters can be directly understood (Component 1)
    \item \textbf{Interpretability:} The model's predictions and behavior can be explained to stakeholders (Components 1-4 combined)
\end{itemize}

A model can be interpretable without being fully transparent (e.g., Random Forest with SHAP explanations), but transparency is a stronger property that typically leads to higher interpretability scores.

\section{Results}

\subsection{Overall Performance Summary}

Table~\ref{tab:overall_summary} presents the aggregate results across all problem domains.

\begin{table}[h]
\centering
\begin{tabular}{lrr}
\toprule
\textbf{Metric} & \textbf{Statistical} & \textbf{ML/DL} \\
\midrule
Avg Performance (R²/Acc) & 0.547 & \textcolor{mlcolor}{0.818} \\
Avg Training Time & \textcolor{statcolor}{0.037s} & 0.763s \\
Efficiency (Performance/Cost) & \textcolor{accentcolor}{High} & Medium \\
Interpretability & \textcolor{accentcolor}{88.3/100} & 37.4/100 \\
\bottomrule
\end{tabular}
\caption{Overall Performance Summary}
\label{tab:overall_summary}
\end{table}

\textbf{Key Observations:}
\begin{itemize}
    \item ML/DL models achieve 1.5x higher average performance (0.818 vs 0.547)
    \item Statistical models are 20.6x faster in training (0.037s vs 0.763s)
    \item Statistical models are 2.4x more interpretable (88.3/100 vs 37.4/100)
    \item Efficiency advantage strongly favors statistical models
\end{itemize}

\subsection{Regression Analysis Results}

Table~\ref{tab:regression_results} shows detailed regression results.

\begin{table}[h]
\centering
\begin{tabular}{lcccc}
\toprule
\textbf{Dataset} & \textbf{OLS R²} & \textbf{ML LR R²} & \textbf{OLS Time} & \textbf{ML Time} \\
\midrule
California Housing & 0.575 & 0.575 & 0.0001s & 0.0005s \\
Diabetes & 0.453 & 0.453 & 0.0001s & 0.0005s \\
Wine Quality & 0.787 & 0.787 & 0.0001s & 0.0005s \\
Energy Consumption & 0.811 & \textcolor{mlcolor}{0.970} & 0.0001s & 0.0005s \\
\bottomrule
\end{tabular}
\caption{Regression Analysis Results}
\label{tab:regression_results}
\end{table}

\textbf{Findings:}
\begin{itemize}
    \item \textbf{Linear Datasets:} OLS and ML Linear Regression achieve identical performance (California Housing, Diabetes, Wine Quality)
    \item \textbf{Complex Patterns:} ML models excel on Energy Consumption dataset (0.970 vs 0.811)
    \item \textbf{Speed Advantage:} OLS is 5x faster (0.0001s vs 0.0005s)
    \item \textbf{Statistical Inference:} Only OLS provides p-values, confidence intervals, and model diagnostics
\end{itemize}

\subsection{Classification Analysis Results}

Table~\ref{tab:classification_results} presents classification performance.

\begin{table}[h]
\centering
\begin{tabular}{lcccc}
\toprule
\textbf{Dataset} & \textbf{Logistic Acc.} & \textbf{RF Acc.} & \textbf{Logistic Time} & \textbf{RF Time} \\
\midrule
Iris & \textcolor{accentcolor}{1.000} & \textcolor{accentcolor}{1.000} & 0.002s & 0.15s \\
Wine Classification & 0.956 & \textcolor{mlcolor}{0.982} & 0.002s & 0.15s \\
Breast Cancer & 0.956 & \textcolor{mlcolor}{0.965} & 0.002s & 0.15s \\
Credit Default & \textcolor{accentcolor}{0.999} & \textcolor{accentcolor}{0.999} & 0.002s & 0.15s \\
\bottomrule
\end{tabular}
\caption{Classification Analysis Results}
\label{tab:classification_results}
\end{table}

\textbf{Findings:}
\begin{itemize}
    \item \textbf{Linearly Separable:} Both approaches achieve perfect accuracy (Iris, Credit Default)
    \item \textbf{Complex Boundaries:} Random Forest superior on Wine Classification and Breast Cancer
    \item \textbf{Speed:} Logistic Regression 75x faster (0.002s vs 0.15s)
    \item \textbf{Interpretability:} Logistic Regression provides odds ratios and p-values
\end{itemize}

\subsection{Factorial Analysis Results}

Table~\ref{tab:factorial_results} shows ANOVA vs SVM comparison.

\begin{table}[h]
\centering
\begin{tabular}{lcccc}
\toprule
\textbf{Dataset} & \textbf{ANOVA R²} & \textbf{SVM R²} & \textbf{ANOVA Time} & \textbf{SVM Time} \\
\midrule
Treatment Effects & 0.417 & \textcolor{mlcolor}{0.513} & 0.0005s & 2.91s \\
Educational Outcomes & -0.002 & \textcolor{mlcolor}{0.284} & 0.001s & 0.19s \\
Agricultural Experiment & 0.464 & \textcolor{mlcolor}{0.745} & 0.001s & 0.23s \\
Drug Efficacy Study & 0.728 & \textcolor{mlcolor}{0.880} & 0.001s & 0.25s \\
\bottomrule
\end{tabular}
\caption{Factorial Analysis Results (ANOVA vs SVM)}
\label{tab:factorial_results}
\end{table}

\textbf{Findings:}
\begin{itemize}
    \item \textbf{Prediction Accuracy:} SVM consistently outperforms ANOVA (average 0.606 vs 0.402)
    \item \textbf{Speed:} ANOVA is 100-1000x faster
    \item \textbf{Interpretability:} ANOVA provides F-statistics, p-values, and interaction effects
    \item \textbf{Use Case:} ANOVA for hypothesis testing, SVM for prediction
\end{itemize}

\subsection{Time Series Analysis Results}

Table~\ref{tab:timeseries_results} presents ARIMA vs LSTM comparison.

\begin{table}[h]
\centering
\begin{tabular}{lcccc}
\toprule
\textbf{Dataset} & \textbf{ARIMA R²} & \textbf{LSTM R²} & \textbf{ARIMA Time} & \textbf{LSTM Time} \\
\midrule
Energy Consumption & 0.450 & \textcolor{dlcolor}{0.890} & 0.18s & 2.5s \\
Temperature & \textcolor{statcolor}{0.972} & 0.965 & 0.18s & 2.5s \\
Sales Forecasting & \textcolor{statcolor}{0.972} & 0.969 & 0.18s & 2.5s \\
\bottomrule
\end{tabular}
\caption{Time Series Analysis Results (ARIMA vs LSTM)}
\label{tab:timeseries_results}
\end{table}

\textbf{Findings:}
\begin{itemize}
    \item \textbf{Complex Patterns:} LSTM excels on Energy Consumption (0.890 vs 0.450)
    \item \textbf{Stationary Data:} ARIMA performs excellently on Temperature and Sales (0.972)
    \item \textbf{Speed:} ARIMA is 14x faster (0.18s vs 2.5s)
    \item \textbf{Interpretability:} ARIMA provides interpretable coefficients and diagnostics
\end{itemize}

\subsection{Performance vs Cost Analysis}

Figure~\ref{fig:cost_performance} visualizes the fundamental trade-off between model performance and training cost.

\begin{figure}[h]
\centering
\includegraphics[width=0.8\textwidth]{CostVsPerformance.png}
\caption{Performance vs Training Cost Trade-off. Statistical models (light blue) cluster in the low-cost, moderate-performance region. ML models (purple) show moderate cost with higher performance. Deep Learning models (orange) achieve highest performance at highest cost. The ideal solution (top-left) represents high performance with low cost.}
\label{fig:cost_performance}
\end{figure}

\textbf{Key Observations:}
\begin{itemize}
    \item \textcolor{statcolor}{Statistical models} occupy the low-cost, moderate-performance quadrant
    \item \textcolor{mlcolor}{ML models} demonstrate better performance-cost balance
    \item \textcolor{dlcolor}{Deep Learning} achieves highest performance but at highest cost
    \item \textbf{Ideal Solution:} Top-left corner (high performance, low cost) - target for future research
\end{itemize}

\subsection{Efficiency Analysis}

Table~\ref{tab:efficiency} ranks models by efficiency (Performance/Cost ratio).

\begin{table}[h]
\centering
\begin{tabular}{lrr}
\toprule
\textbf{Ranking} & \textbf{Model} & \textbf{Efficiency} \\
\midrule
1 & OLS Regression (Stat) & \textcolor{accentcolor}{6435.64} \\
2 & ANOVA (Stat) & \textcolor{accentcolor}{5445.54} \\
3 & Linear Regression (ML) & 1297.41 \\
4 & Logistic Regression (Stat) & 439.78 \\
5 & Poisson Regression (Stat) & 349.65 \\
\ldots & \ldots & \ldots \\
Last & RNN/LSTM (DL) & 0.36 \\
\bottomrule
\end{tabular}
\caption{Efficiency Rankings (Performance/Cost)}
\label{tab:efficiency}
\end{table}

\textbf{Insights:}
\begin{itemize}
    \item Statistical models dominate efficiency rankings (top 4 positions)
    \item OLS Regression achieves 17,800x higher efficiency than LSTM
    \item Efficiency advantage makes statistical models ideal for real-time applications
\end{itemize}

\subsection{Winners by Category}

Table~\ref{tab:winners} identifies best-performing models across different criteria.

\begin{table}[h]
\centering
\begin{tabular}{lcc}
\toprule
\textbf{Category} & \textbf{Winner} & \textbf{Performance} \\
\midrule
Best Overall Accuracy & Random Forest (ML) & 0.920 \\
Fastest Training & OLS Regression (Stat) & 0.0001s \\
Most Interpretable & OLS Regression (Stat) & 100/100 \\
Best Efficiency & OLS Regression (Stat) & 6435.64 \\
Best Time Series (Simple) & ARIMA (Stat) & R² = 0.972 \\
Best Time Series (Complex) & LSTM (DL) & R² = 0.890 \\
\bottomrule
\end{tabular}
\caption{Winners by Category}
\label{tab:winners}
\end{table}

\section{Discussion}

\subsection{The Fundamental Trade-offs}

Our results confirm the existence of three fundamental trade-offs in model selection:

\subsubsection{Performance vs Interpretability}

The traditional inverse relationship is clearly evident (Figure~\ref{fig:interpretability_tradeoff}). Statistical models achieve high interpretability (88.3/100) but moderate performance (0.547), while ML/DL models achieve high performance (0.818) but lower interpretability (37.4/100).

However, the field of \textbf{Interpretable ML} is actively working to bridge this gap. Research on Transparent AI \cite{transparent_ai_2025} demonstrates that interpretability can be integrated into ML workflows without sacrificing all performance gains.

\subsubsection{Performance vs Cost}

Figure~\ref{fig:cost_performance} illustrates the performance-cost trade-off. Statistical models provide excellent cost efficiency, while ML/DL models offer superior performance at higher computational cost.

\subsubsection{Interpretability vs Cost}

Interestingly, interpretability and cost are not necessarily in conflict. Statistical models achieve both high interpretability and low cost, making them ideal for applications requiring transparency and real-time performance.

\subsection{When to Choose Statistical Models}

Based on our findings, statistical models are recommended when:

\begin{enumerate}
    \item \textbf{Interpretability is Critical:} Need for p-values, confidence intervals, and statistical inference (e.g., regulatory compliance, scientific research)
    \item \textbf{Regulatory Compliance:} Transparent decision-making required (healthcare, finance, public policy)
    \item \textbf{Hypothesis Testing:} Understanding causal relationships and testing theories
    \item \textbf{Small to Medium Datasets:} Optimal performance on limited data (statistical models often outperform ML on small datasets)
    \item \textbf{Real-time Applications:} Fast training and prediction needed (20-100x faster)
    \item \textbf{Research and Development:} Need to understand underlying relationships
\end{enumerate}

\textbf{Best Statistical Models:} OLS Regression, Logistic Regression, ANOVA, ARIMA, Cox Regression

\subsection{When to Choose ML/DL Models}

ML/DL models are recommended when:

\begin{enumerate}
    \item \textbf{Maximum Accuracy is Priority:} Need best predictive performance (1.5x higher on average)
    \item \textbf{Large, Complex Datasets:} High-dimensional or non-linear patterns
    \item \textbf{Production Systems:} Computational resources available for training
    \item \textbf{Pattern Discovery:} Exploratory analysis on complex data
    \item \textbf{Non-linear Relationships:} When linear assumptions fail
    \item \textbf{Sequential/Temporal Data:} LSTM/RNN for time series with complex patterns
\end{enumerate}

\textbf{Best ML/DL Models:} Random Forest, SVM, LSTM/RNN, Gradient Boosting

\subsection{Implications for Transparent AI}

Our findings align with the principles outlined in Transparent AI research \cite{transparent_ai_2025}:

\begin{itemize}
    \item \textbf{Early Integration:} Model selection should consider interpretability requirements from the start
    \item \textbf{Stakeholder Needs:} Different audiences require different explanation types - statistical models excel for domain experts, while ML models may require post-hoc explanation methods
    \item \textbf{Regulatory Alignment:} Statistical models provide natural compliance with transparency requirements
    \item \textbf{Hybrid Approaches:} Future research should explore combining statistical inference with ML performance
\end{itemize}

\subsection{Limitations and Future Work}

\subsubsection{Limitations}

\begin{itemize}
    \item \textbf{Dataset Selection:} Limited to 13 datasets - broader evaluation needed
    \item \textbf{Hyperparameter Tuning:} ML models may benefit from more extensive tuning
    \item \textbf{Interpretability Metrics:} Our interpretability score is composite - domain-specific evaluation needed
    \item \textbf{Post-hoc Explanations:} We did not evaluate SHAP, LIME, or other explanation methods for ML models
\end{itemize}

\subsubsection{Future Research Directions}

\begin{enumerate}
    \item \textbf{Hybrid Models:} Develop models combining statistical inference with ML performance
    \item \textbf{Interpretable ML:} Evaluate inherently interpretable ML models (EBMs, GAMs, NAMs)
    \item \textbf{Post-hoc Explanations:} Comprehensive evaluation of SHAP, LIME, and other methods
    \item \textbf{Domain-Specific Studies:} Healthcare, finance, and other high-stakes applications
    \item \textbf{Regulatory Frameworks:} Alignment with EU AI Act, Canada AIDA, and other regulations
\end{enumerate}

\section{Conclusions}

\subsection{Key Takeaways}

\begin{enumerate}
    \item \textbf{No Universal Winner:} Both approaches have distinct advantages - the choice depends on specific requirements
    
    \item \textbf{Statistical Models Excel In:}
    \begin{itemize}
        \item Interpretability (88.3/100 vs 37.4/100) - 2.4x advantage
        \item Speed (20-100x faster training)
        \item Efficiency (6435.64 vs 0.36 efficiency scores)
        \item Regulatory compliance and hypothesis testing
    \end{itemize}
    
    \item \textbf{ML/DL Models Excel In:}
    \begin{itemize}
        \item Complex dataset performance (0.818 vs 0.547 avg) - 1.5x advantage
        \item Scalability to large data
        \item Capturing non-linear patterns
        \item Handling complex temporal dependencies
    \end{itemize}
    
    \item \textbf{The Optimal Choice Depends On:}
    \begin{itemize}
        \item Primary requirement (accuracy vs interpretability vs speed)
        \item Dataset characteristics (size, complexity, linearity)
        \item Computational constraints
        \item Regulatory requirements
    \end{itemize}
\end{enumerate}

\subsection{Strategic Recommendations}

Rather than choosing one approach over the other, organizations should:

\begin{enumerate}
    \item \textbf{Develop Expertise in Both:} Build capabilities in both statistical and ML methodologies
    \item \textbf{Context-Aware Selection:} Choose models based on specific project requirements
    \item \textbf{Hybrid Approaches:} Explore combining statistical inference with ML performance
    \item \textbf{Interpretability Integration:} Embed interpretability considerations from initial design
    \item \textbf{Regulatory Alignment:} Ensure model selection aligns with compliance requirements
\end{enumerate}

\subsection{Contribution to Transparent AI}

This study contributes to the growing body of research on Transparent AI by:

\begin{itemize}
    \item Providing empirical evidence on performance-interpretability trade-offs
    \item Quantifying the efficiency advantages of statistical models
    \item Demonstrating when each approach is most appropriate
    \item Connecting findings to regulatory requirements and stakeholder needs
    \item Identifying opportunities for hybrid approaches
\end{itemize}

As AI systems increasingly inform critical decisions, the ability to balance performance, cost, and interpretability becomes essential. This study provides evidence-based guidance for practitioners navigating these complex trade-offs in the age of Transparent AI.

\section{Acknowledgments}

We thank the Master Paris-Saclay Applied Statistics program for supporting this research. We also acknowledge the open-source communities developing statistical and machine learning tools that made this comparative study possible.

\bibliographystyle{plain}
\begin{thebibliography}{9}

\bibitem{transparent_ai_2025}
Ramachandram, D., Joshi, H., Zhu, J., Gandhi, D., Hartman, L., \& Raval, A. (2025).
\textit{Transparent AI: The Case for Interpretability and Explainability}.
arXiv:2507.23535v1 [cs.LG]

\bibitem{rudin2019}
Rudin, C., \& Radin, J. (2019).
Why Are We Using Black Box Models in AI When We Don't Need To? A Lesson From An Explainable AI Competition.
\textit{Harvard Data Science Review}.

\bibitem{arrieta2020}
Arrieta, A. B., et al. (2020).
Explainable Artificial Intelligence (XAI): Concepts, Taxonomies, Opportunities and Challenges.
\textit{Information Fusion}, 58, 82-115.

\bibitem{hastie2009}
Hastie, T., Tibshirani, R., \& Friedman, J. (2009).
\textit{The Elements of Statistical Learning}.
Springer.

\bibitem{james2013}
James, G., Witten, D., Hastie, T., \& Tibshirani, R. (2013).
\textit{An Introduction to Statistical Learning}.
Springer.

\bibitem{murphy2012}
Murphy, K. P. (2012).
\textit{Machine Learning: A Probabilistic Perspective}.
MIT Press.

\bibitem{cox1972}
Cox, D. R. (1972).
Regression Models and Life-Tables.
\textit{Journal of the Royal Statistical Society}, Series B, 34(2), 187-220.

\bibitem{kaplan1958}
Kaplan, E. L., \& Meier, P. (1958).
Nonparametric Estimation from Incomplete Observations.
\textit{Journal of the American Statistical Association}, 53(282), 457-481.

\end{thebibliography}

\end{document}

